% Teaching manuscript. Edit this TeX source; Markdown is compatibility navigation only.

\StudyChapter{CS5222-ResearchReportGuide}{Research Report Guide｜把一个真实问题解释清楚}
\StudyTarget{CS5222-ResearchReportGuide__research-report-guide}


\StudyRepoRef{CS5222/course-notes/README.md}{}{课程目录} · \StudyRef{CS5222-Chapter01}{root}{Chapter1} · \StudyRef{CS5222-Chapter02}{root}{Chapter2} · \href{https://crazyshout.github.io/micro-course/notices.html}{课程信息栏}


这份指南帮助你选题、解释机制、核查资料和整理反思。原任务要求研究你在日常生活、新闻或工作中遇到的新网络技术，联系教材Chapter1–7，排除Chapter8安全主题；用自己的解释和图说明机制，并反思学习过程。格式、页数、评分与提交要求以Canvas中的当前原任务与Rubric为准。


按选题补基础：涉及码率/吞吐就回\StudyRef{Foundations-NetworkBasics}{units}{网络单位}，涉及状态/时间就回\StudyRef{Foundations-NetworkBasics}{messages}{报文与时间轴}。没有实际卡点便直接开始。


\begingroup
\small\setlength{\tabcolsep}{3pt}
\renewcommand{\arraystretch}{1.22}
\begin{longtable}{@{}>{\raggedright\arraybackslash}p{\dimexpr0.212307\linewidth-1.50000\tabcolsep\relax}>{\raggedright\arraybackslash}p{\dimexpr0.186423\linewidth-1.50000\tabcolsep\relax}>{\raggedright\arraybackslash}p{\dimexpr0.235831\linewidth-1.50000\tabcolsep\relax}>{\raggedright\arraybackslash}p{\dimexpr0.365438\linewidth-1.50000\tabcolsep\relax}@{}}
\toprule
能力 & 优先级 & 难度与卡点 & 达标表现；依据 \\
\midrule\endfirsthead
能力 & 优先级 & 难度与卡点 & 达标表现；依据 \\
\midrule\endhead
\bottomrule\endfoot
明确问题、解释机制 & 核心必会 & 中至高；不能只有名词 & 不看原文也能画出数据/控制流；原任务A \\
来源核查与边界 & 核心必会 & 中；资料版本/适用场景 & 每个关键结论有对应证据；原任务p2 \\
图文与英语表达 & 核心必会 & 中；精简与可读 & 图能帮助解释一个具体难点；原Rubric \\
真实反思 & 核心必会 & 中；具体变化而非套话 & 记录自己确实做过的核查与困难；原任务B \\
复杂数学/未用机制 & 按需拓展 & 可能高 & 只有支撑论点时补；整理者建议，不暗示不考 \\
\end{longtable}
\endgroup



\StudyHeading{1}{1. 从实际问题确定研究范围}
\StudyTarget{CS5222-ResearchReportGuide__1}


先记录你\textbf{确实遇到}的情况：某个应用在带宽变化时切画质、同一内容从不同地点下载速度不同、多人共享网络时服务突然卡顿。题目应围绕一个能解释的机制问题，例如“客户端怎样在吞吐波动时选择下一视频块的编码版本”。


用三个问题缩小范围：这个问题属于网络吗？能对应教材哪一层/哪章？在本次篇幅内能用一幅图、一个小场景说明主要机制吗？“所有5G技术”往往太宽；“某种特定机制在什么条件下改善什么问题”更容易论证。选题的新意应落在具体变化、部署或问题上。若选DASH等课内基础，需说明你研究的具体新变化、部署或问题，不能只复述课件。


\textbf{English framing (整理表达):} \emph{This report examines how [a specific mechanism] addresses [a concrete networking problem] under [a defined setting].} 方括号中填写你的具体问题、机制和调查情境。



\StudyHeading{1}{2. 把问题、机制、效果连成一条可核查的解释}
\StudyTarget{CS5222-ResearchReportGuide__2}


先写旧方法遇到什么限制，再写新机制改了哪一步，最后讨论收益与代价。不要“协议很先进，所以性能很好”这样跳过中间因果链。


以分块视频为\textbf{教学例}：播放每秒消耗编码数据 \ensuremath{\to} 网络吞吐会变动 \ensuremath{\to} 缓冲可能耗尽 \ensuremath{\to} 客户端按估计吞吐/缓冲选择下一块 \ensuremath{\to} 可能减少停顿，但画质也会变化。这里每条箭头都能被课程概念支持，而不是把“自适应”三个字当解释。


\begin{StudyBlock}{完整示范 / Worked example}{下载时间与播放缓冲}
4秒片段编码2Mbps，数据量8Mbit；若稳定吞吐4Mbps，下载2秒。在初始缓冲足够、下载中同时播放、整块到齐后加入缓冲且忽略请求与解码开销时，缓冲净增加4\ensuremath{-}2=2秒。 / A four-second segment encoded at 2 Mbps contains 8 Mbit. At a constant 4-Mbps download rate it arrives in two seconds. With sufficient initial buffer, playback continuing during download, whole-chunk availability on arrival and no request or decoding overhead, playable buffer grows by two seconds.
\end{StudyBlock}


\textbf{变式 / Transfer:} 一个4秒视频块按2 Mbps编码，下载吞吐稳定为1 Mbps，开始时有3秒可播放缓冲。下载中继续播放，整块到齐才可播放，忽略其他开销。下载需要多久，会停顿多久？ / A four-second chunk encoded at 2 Mbps downloads at a constant 1 Mbps, starting with three seconds of playable buffer. Playback continues during download; the chunk becomes playable only when fully received. Ignore other overhead. Find download and stall times.



\begin{StudySolution}{答案 / Answer}


下载8秒，已有缓冲只能撑3秒，理想模型中会停顿约5秒，等块到后才续播。 / Download takes eight seconds; under the simplified model, playback stalls for about five seconds.


\end{StudySolution}


播放器策略和实际吞吐变化会改变停顿时间。



\StudyHeading{1}{3. 找来源时，为每一句重要话准备证据}
\StudyTarget{CS5222-ResearchReportGuide__3}


优先找协议标准、技术提供方文档、原始研究论文或可复核的测量报告。二手讲解和AI适合帮助形成问题，但需要回到原资料核对。阅读来源时，找到直接支持该结论的段落或数据。


\begingroup
\small\setlength{\tabcolsep}{3pt}
\renewcommand{\arraystretch}{1.22}
\begin{longtable}{@{}>{\raggedright\arraybackslash}p{\dimexpr0.280642\linewidth-1.33333\tabcolsep\relax}>{\raggedright\arraybackslash}p{\dimexpr0.332215\linewidth-1.33333\tabcolsep\relax}>{\raggedright\arraybackslash}p{\dimexpr0.387143\linewidth-1.33333\tabcolsep\relax}@{}}
\toprule
写作中的主张 & 应找的证据 & 必须检查 \\
\midrule\endfirsthead
写作中的主张 & 应找的证据 & 必须检查 \\
\midrule\endhead
\bottomrule\endfoot
协议怎样工作 & 规范/原设计文档 & 版本、消息顺序、可选项、条件 \\
某产品使用此机制 & 官方当前产品文档 & 原生/浏览器/企业部署是否相同 \\
性能提高多少 & 原实验/可重复测量 & 对比方法、数据、指标、环境 \\
有哪些限制 & 规范限制、实验失败或设计分析 & 不把推断写成已测事实 \\
为什么和课程有关 & 教材/当前课件定位 & 是当前已学，还是教材后续章节 \\
\end{longtable}
\endgroup


可维护一张很小的工作表，放在你自己的研究草稿里：\StudyInline{Claim | Source/\allowbreak{}section | Version/\allowbreak{}date | Supports what | Does not support what}。用它检查每条主张与来源是否对应。


如引用AI输出，原任务要求说明使用并自行负责核查；图、引用也要查，不能只给一个模型名字就当技术来源。对论文核对题名、作者、年份、DOI/官方页面，打开原文确认所引结论和实验设置。查不到原文的来源标为待核查，确认后再引用。



\StudyHeading{1}{4. 用图说明交互、结构或权衡}
\StudyTarget{CS5222-ResearchReportGuide__4}


选择一种与你的核心问题匹配的图：



\begin{itemize}

\item 交互问题画时序图，标清参与者、请求、响应、时间方向；

\item 架构问题画节点与数据/控制流，说明图中箭头是什么；

\item 权衡问题画有单位的曲线，说明数据来自实测、论文还是教学计算。

\end{itemize}


画完问自己：读者通过这张图能更容易答出哪个问题？节点之间的箭头应说明数据从哪里来、到哪里去，以及每一步做什么。引用原图标来源，改编图标“Adapted from…”，自己计算的图保存脚本和输入。AI图同样要署明并逐条检查技术结构，不能因为好看就默认网络方向正确。



\StudyHeading{1}{5. 英文报告怎样组织才连贯}
\StudyTarget{CS5222-ResearchReportGuide__5}


建议采用内容驱动的小标题，而非每段中英对照：英文完成正式论证，中文草稿帮助自己先讲懂。术语第一次出现给全称与缩写，随后保持一致。可以按以下结构组织：



\begin{enumerate}[start=1]

\item \textbf{Problem and context}：你实际怎样接触主题？问题是什么、研究范围到哪里？

\item \textbf{Background from the course}：挑真正用到的1–3个概念，解释它们为何必要。

\item \textbf{How the mechanism works}：沿一个具体请求/数据流走完步骤，配关键图。

\item \textbf{Benefits, limits and evidence}：哪项收益有证据，代价是什么，什么条件下不适用？

\item \textbf{Discussion}：你基于证据形成的判断、未解决问题与合理后续方向。

\end{enumerate}


每段先给主张，再解释机制，再放证据与条件。比如 \emph{The mechanism reduces repeated origin fetches when reusable content is available in cache} 比 \emph{It greatly improves everything} 更可核查。保留影响结论的条件和不确定性。


\textbf{英文自测 / English self-check:} \emph{What changes, why does it help, and under what conditions might it not help? / 改变了什么，为何有效，何时可能无效？} 若不能脱稿回答，正文还需要补机制，而非继续增加术语。



\StudyHeading{1}{6. Reflection log｜记录理解怎样改变}
\StudyTarget{CS5222-ResearchReportGuide__6-reflection-log}


原任务的五个问题，可以分别对应一条具体记录：


\begingroup
\small\setlength{\tabcolsep}{3pt}
\renewcommand{\arraystretch}{1.22}
\begin{longtable}{@{}>{\raggedright\arraybackslash}p{\dimexpr0.611844\linewidth-1.00000\tabcolsep\relax}>{\raggedright\arraybackslash}p{\dimexpr0.388156\linewidth-1.00000\tabcolsep\relax}@{}}
\toprule
原问题 & 真实记录的抓手 \\
\midrule\endfirsthead
原问题 & 真实记录的抓手 \\
\midrule\endhead
\bottomrule\endfoot
What did you learn that you did not know before? & 原来怎么理解，哪个具体事实改变了它 \\
What was the most challenging part? & 一个真正卡住的机制/推导/资料冲突 \\
How did GenAI help or hinder? & 它帮你拆了哪一步，或给过什么误导 \\
How did you verify AI-generated information? & 打开了哪个原来源、比对哪句、做了什么检查 \\
What would you do differently next time? & 一个具体可实行的改进 \\
\end{longtable}
\endgroup


可从现在开始记两三条带日期的研究记录，报告结束时据实整理。写清你做过什么、哪里卡住、怎样查证，以及理解发生了什么变化；工具使用情况按实际记录。



\StudyHeading{1}{7. 交付前的自查与来源覆盖}
\StudyTarget{CS5222-ResearchReportGuide__7}


读者应能复述你的问题与机制，认出它对应课程哪部分；一张图能独立读懂；数字都有单位、来源和条件；观点与已验证事实分开。英语表达、引用、AI使用说明、格式与评分要求逐项回查原任务及课程信息。


\begingroup
\small\setlength{\tabcolsep}{3pt}
\renewcommand{\arraystretch}{1.22}
\begin{longtable}{@{}>{\raggedright\arraybackslash}p{\dimexpr0.593176\linewidth-1.00000\tabcolsep\relax}>{\raggedright\arraybackslash}p{\dimexpr0.406824\linewidth-1.00000\tabcolsep\relax}@{}}
\toprule
原任务位置 & 本指南对应 \\
\midrule\endfirsthead
原任务位置 & 本指南对应 \\
\midrule\endhead
\bottomrule\endfoot
p1 Goal、题目范围 & 开篇、§1 \\
p1 Format、提交构成 & 原任务/信息栏链接；不生成正式提交件 \\
p1 A：接触途径、机制、联系、动机/创新/挑战/观点 & §1–2、§5 \\
p1 自己的表达、图与清晰度 & §4–5 \\
p2 GenAI引用、图/参考核查责任 & §3–4 \\
p2 B：五个Reflection问题 & §6逐项对应 \\
p2 Evaluation、p3 Rubric & 开篇能力表、§7；具体权重以原页为准 \\
\end{longtable}
\endgroup
